\documentclass[11pt,a4paper]{article}
\usepackage[utf8]{inputenc}
\usepackage[T1]{fontenc}
\usepackage[spanish,es-nodecimaldot]{babel}
\usepackage[margin=2cm]{geometry}
\usepackage{amsmath,amssymb,booktabs,longtable,tabularx,array}
\usepackage{enumitem,xcolor,hyperref,xurl}
\hypersetup{colorlinks=true,linkcolor=blue!55!black,urlcolor=blue!55!black}
\setlist{nosep,leftmargin=*}
\newcommand{\pend}{\textbf{Pendiente de confirmar}}
\newcommand{\prop}{\textbf{Propuesta}}
\newcommand{\hecho}{\textbf{Hecho verificado}}
\title{Flujo DG2 para una futura simulación de inundación\\\large TerraPulse AI}
\author{Documento de diseño; no implementa ni autoriza una simulación}
\date{Dirección comunicada: DG2; requisitos y versión concreta pendientes}
\begin{document}
\maketitle
\begin{abstract}
El equipo comunicó que el físico eligió DG2 como dirección de trabajo. Este documento presenta esa ruta sin atribuirle un motivo único; la revisión física y numérica deberá documentar el régimen, las hipótesis de aguas someras, la pertinencia de la advección, la calidad de los datos y las pruebas de estabilidad, conservación y convergencia. El alcance acordado se limita a terrenos no urbanos: no se modelarán ciudades ni variables urbanas. Este documento separa el flujo físico del flujo de aplicación, distingue evidencia publicada para LISFLOOD-FP 8.0 de capacidades pendientes del binario que eventualmente se elija y no inventa evento, entradas, parámetros ni resultados.
\end{abstract}
\tableofcontents\newpage

\section{Estado de evidencia y decisión}
\begin{description}
\item[Dirección del equipo.] DG2 es la ruta seleccionada para diseñar la propuesta; no equivale a que un ejecutable, versión, evento o parámetros estén aprobados.
\item[Alcance acordado.] Estudiar inundación sobre terreno no urbano, sin incorporar edificios, calles, drenaje pluvial urbano ni variables de ciudades. La topografía, cauces naturales, rugosidad del terreno y fronteras hidráulicas siguen siendo necesarios. La selección del lugar no está hecha: no se presume que la ROI o su dominio completo estén libres de zonas urbanas.
\item[Justificación de DG2.] Se requieren evidencia de que aceleración advectiva/transitorios importan y que las hipótesis de aguas someras son adecuadas; necesidad de resolver gradientes/frentes mejor que un método de primer orden; benchmarks pertinentes con conservación y convergencia; datos de terreno, cauce y fronteras adecuados; estabilidad y costo medidos en el dominio previsto. DG2 no garantiza por sí mismo mejor exactitud física.
\item[Pendientes.] Evento y régimen; verificación de que dominio y área contribuyente sean compatibles con el alcance no urbano; DEM, datum y cauces; lluvia/caudales y fronteras; infiltración; rugosidad; malla y periodos; versión exacta, manual, binario, opciones, formatos y licencia; benchmarks y tolerancias; observación independiente para validar.
\end{description}

\part{I. Flujo físico}
\section{Hipótesis y ecuaciones del sistema}
\prop\ Representar agua somera promediada verticalmente con profundidad $h(x,y,t)$, velocidad horizontal media $\boldsymbol u=(u,v)$ y elevación del lecho $z_b$. Se supone escala horizontal mayor que profundidad y presión aproximadamente hidrostática. El sistema no resuelve flujo 3D local, turbulencia detallada ni erosión. Este caso no representará edificios, calles, alcantarillado urbano ni barreras construidas: si alguno altera significativamente el flujo hacia o dentro del dominio, el caso no cumple este alcance y debe replantearse, no simularse como si la barrera no existiera. Corresponde comprobar estas hipótesis para el evento, no asumirlas.

Una forma conceptual conservativa de SWE 2D es
\begin{align}
\partial_t h+\nabla\cdot(h\boldsymbol u)&=r-i+s_h, \label{eq:cont}\\
\partial_t(h\boldsymbol u)+\nabla\cdot\left(h\boldsymbol u\otimes\boldsymbol u+\tfrac12gh^2 I\right)&=-gh\nabla z_b-gh\boldsymbol S_f+\boldsymbol a. \label{eq:mom}
\end{align}
$\boldsymbol S_f$ aquí significa pendiente de fricción; en la notación de la fuente LISFLOOD-FP 8.0, $S_f$ puede designar el término de fuerza de fricción con signo, no esta pendiente adimensional. No confundir las dos convenciones. La aceleración/fuente de momento $\boldsymbol a$ solo se incluiría con definición y soporte verificados.

\paragraph{Comprobación dimensional término a término.} Las derivadas temporales dividen entre segundos y las espaciales entre metros horizontales. En continuidad, $\partial_t h$ es profundidad sobre el lecho por tiempo $[\mathrm{m/s}]$; $h\boldsymbol u$ es volumen por tiempo y por ancho perpendicular al flujo $[\mathrm{m^2/s}]$, de modo que $\nabla\cdot(h\boldsymbol u)$ divide por distancia horizontal y resulta $[\mathrm{m/s}]$. La lluvia $r$, la infiltración $i$ y la fuente neta $s_h$ son cada una volumen por área horizontal de planta y por tiempo $[\mathrm{m^3/(m^2s)}=\mathrm{m/s}]$; sus signos no cambian la dimensión. Por tanto, ambos lados de \eqref{eq:cont} están en m/s.

En cada componente horizontal de momento, $h\boldsymbol u$ representa volumen por tiempo y ancho $[\mathrm{m^2/s}]$ y $\partial_t(h\boldsymbol u)$ añade un divisor temporal: $[\mathrm{m^2/s^2}]$. El flujo advectivo $h\boldsymbol u\otimes\boldsymbol u$ es $[\mathrm m(\mathrm{m/s})^2=\mathrm{m^3/s^2}]$; el flujo de presión $gh^2 I/2$ es $[(\mathrm{m/s^2})\mathrm{m^2}=\mathrm{m^3/s^2}]$ (el $1/2$ e $I$ son adimensionales). La divergencia de \emph{cada} flujo divide por distancia horizontal $[1/\mathrm m]$: ambos aportan $[\mathrm{m^2/s^2}]$. El gradiente del lecho $\nabla z_b$ compara elevaciones respecto del mismo datum por distancia horizontal $[\mathrm{m/m}=1]$; $-gh\nabla z_b$ y $-gh\boldsymbol S_f$ son cada uno $[(\mathrm{m/s^2})\mathrm m(1)=\mathrm{m^2/s^2}]$. La fuente $\boldsymbol a$ debe representar aportes de momento por área horizontal y densidad del agua (dimensión $[\mathrm{m^2/s^2}]$), incluida, si procede, la cantidad de movimiento transportada por fuentes/sumideros de masa: no se presume independiente de $r,i,s_h$. Así todos los sumandos de \eqref{eq:mom} tienen m$^2$/s$^2$; esta igualdad dimensional no prueba la validez física de las fuentes ni su implementación en el solver.

\section{Fuentes, fricción y operadores: glosario dimensional}
\begin{longtable}{p{.17\linewidth}p{.18\linewidth}p{.57\linewidth}}
\toprule Símbolo & Unidad & Significado y comprobación \\\midrule
$r$ & m/s & Lluvia que llega a la superficie como profundidad añadida por unidad de tiempo sobre el área modelada; si se resta $i$ por separado, $r$ no debe incorporar ya esa misma pérdida. Una lámina de lluvia $r\Delta t$ tiene metros. Que un binario admita lluvia y su formato concreto se debe verificar; no se infiere de la ecuación. \\
$i$ & m/s & Infiltración como profundidad retirada del agua superficial por unidad de tiempo; aparece restando porque reduce almacenamiento superficial. No representa automáticamente pérdidas subterráneas complejas y no debe fijarse tácitamente en cero. El soporte del ejecutable está pendiente. \\
$s_h$ & m/s & Otras fuentes netas areales: signo positivo añade agua y negativo la retira. Identificar proceso, área y convención; no es una entrada genérica garantizada del solver. \\
$Q$, $q$, $B$ & m$^3$/s, m$^2$/s, m & $Q$ es el volumen que cruza una sección transversal o frontera por tiempo (m$^3$/s); \emph{en este documento}, $q$ es caudal por ancho horizontal de esa sección (m$^3$/(s\,m)). Solo si la distribución es uniforme sobre un ancho $B$ (m), $q=Q/B$; en general $Q=\int_B q\,\mathrm d\ell$. No confundir el ancho con un área. \\
$A_c$, $Q/A_c$ & m$^2$, m/s & $A_c$ es el área transversal \emph{mojada} (m$^2$); $Q/A_c$ es velocidad media del flujo a través de esa sección (m/s), no lluvia ni fuente superficial. \\
$A_p$, $Q/A_p$ & m$^2$, m/s & $A_p$ es área horizontal en planta (m$^2$) sobre la que se \emph{reparte} un volumen de entrada. Solo en esa hipótesis $Q/A_p$ es tasa equivalente de altura de lámina por tiempo: $(Q/A_p)A_p\Delta t=Q\Delta t$ (m$^3$). Una entrada por frontera no se convierte automáticamente en fuente areal ni se cuenta dos veces; la interfaz del solver queda sin verificar. \\
$n$ & s/m$^{1/3}$ & Rugosidad de Manning en unidades SI usuales; fuente, asignación espacial y convención pendientes. \\
$S_{f,x},S_{f,y}$ & 1 & Componentes adimensionales de pendiente de fricción. Una forma de Manning ilustrativa para $h>0$ es $S_{f,x}=n^2u\sqrt{u^2+v^2}/h^{4/3}$ y análogamente para $y$. Unidades: $(\mathrm{s^2/m^{2/3}})(\mathrm{m^2/s^2})/(\mathrm{m^{4/3}})=1$. No implica que sea la forma discreta del ejecutable. \\
$gh\boldsymbol S_f$ & m$^2$/s$^2$ & $(\mathrm{m/s^2})(\mathrm m)(1)$, unidad de cada término de momento. \\
$I$ & 2$\times$2, adimensional & Matriz identidad; $gh^2I$ tiene m$^3$/s$^2$, como el flujo advectivo tensorial antes de aplicar divergencia. \\
$\nabla$ & 1/m antes de actuar & Operador espacial. Sobre $h$ produce $\nabla h$ adimensional; sobre $z_b$ produce pendiente adimensional; divergencia de $h\boldsymbol u$ da m/s; divergencia del tensor de flujo de momento da m$^2$/s$^2$. \\
$\otimes$ & depende de factores & Producto tensorial: $\boldsymbol u\otimes\boldsymbol u$ tiene m$^2$/s$^2$, y al multiplicar por $h$ da m$^3$/s$^2$. \\
\bottomrule
\end{longtable}
\paragraph{Coordenadas y referencias restantes.} $x$ e $y$ son coordenadas cartesianas horizontales de una proyección métrica; sus diferencias y las longitudes $B$, $\mathrm d\ell$ son metros. $t$ y $\Delta t$ son segundos transcurridos respecto a un origen temporal declarado (fechas absolutas en UTC cuando corresponda). $h$ es altura vertical local de agua sobre el lecho (m); $z_b$ es elevación vertical del lecho sobre un datum común (m). Si se usa $\eta=z_b+h$, es elevación de la superficie libre sobre ese mismo datum (m), no profundidad respecto del lecho. $u,v$ y $\boldsymbol u$ son velocidades horizontales medias sobre la columna de agua: desplazamiento en $x,y$ por tiempo (m/s); $g$ es aceleración gravitatoria (m/s$^2$), cambio de velocidad por tiempo. $\boldsymbol S_f$ expresa pérdida de carga vertical equivalente por distancia horizontal (m/m, adimensional); $n$ expresa rugosidad de Manning en la convención SI dada; $I$ es identidad adimensional. $\nabla$ deriva respecto de distancia horizontal (m$^{-1}$ antes de actuar), y $\otimes$ multiplica componentes del vector velocidad, de modo que $\boldsymbol u\otimes\boldsymbol u$ tiene m$^2$/s$^2$. $\boldsymbol a$ es fuente de momento por área horizontal y densidad de agua, integrada en la profundidad (m$^2$/s$^2$); su valor y tratamiento de intercambios de masa quedan por definir. El glosario distingue $A_c$ transversal de $A_p$ en planta: la notación abreviada $Q/A$ es ambigua sin indicar cuál de las dos áreas se usa. \textbf{Correspondencia con la guía del proyecto:} en \path{docs/guia-simulacion-inundacion.md}, §4.2 y §4.6, el símbolo aislado $q$ significa \emph{otras fuentes areales} (m/s), es decir, el $s_h$ de \eqref{eq:cont}; \emph{no} significa el caudal por ancho $q$ (m$^2$/s) definido aquí. En esa misma guía $q_x,q_y$ en las caras de celda designan flujos de volumen por tiempo y ancho (m$^2$/s), compatibles con las componentes de $h\boldsymbol u$; se distinguen del $q$ aislado de fuente. Las ecuaciones de ambos documentos solo se comparan después de esta traducción explícita, nunca sustituyendo $q$ aislado por un flujo de cara.

En celdas secas ($h=0$), la expresión de Manning no se evalúa literalmente: tratamiento, umbrales y regularización dependen de la versión y deben verificarse. Ningún símbolo de fuente demuestra por sí solo que el binario acepte una opción de entrada correspondiente.

\section{Entradas físicas y fronteras}
\begin{itemize}
\item \textbf{Estado inicial:} $h(x,y,0)$ y, si se requiere, velocidad/caudal iniciales; procedencia, hora de referencia y agua permanente explícitas.
\item \textbf{Dominio y ROI:} la ROI es el área que selecciona o consulta la persona usuaria; el dominio hidráulico puede extenderse aguas arriba/abajo e incluir área contribuyente y fronteras relevantes. Verificar ambas áreas como no urbanas para este caso; si el flujo depende de infraestructura urbana omitida, rechazar el caso o cambiar explícitamente el alcance. No recortar automáticamente el dominio a la pantalla.
\item \textbf{Terreno:} DEM con CRS métrico, resolución, nodata, cobertura, fuente, licencia y datum vertical; niveles y DEM deben compartir datum. Verificar cauces naturales y, si importan, su geometría o batimetría; un DEM superficial no garantiza representar el lecho sumergido. No preparar capas de edificios, calles ni drenaje urbano. Si hay estructuras que gobiernan el flujo, documentar esa limitación y no aceptar el caso bajo este alcance simplificado.
\item \textbf{Forzamientos:} lluvia temporal/espacial, caudales o niveles de frontera con unidades, periodo, zona contribuyente, referencia vertical y procedencia. Distinguir entrada por frontera de fuente areal para evitar doble conteo.
\item \textbf{Pérdidas y parámetros:} especificar infiltración y su método/datos, o justificar su omisión y sensibilidad; mapa y fuente de Manning para terreno no urbano, sin categorías urbanas; gravedad, malla, paso y duración tras revisión experta, sin inventar valores.
\end{itemize}

\section{Método numérico, estabilidad y aceptación científica}
\hecho\ Shaw et al. describen LISFLOOD-FP 8.0 con formulación SWE conservativa y un método DG2: representación planar por elemento, integración temporal Runge--Kutta de segundo orden y flujos HLL en la publicación. El artículo reporta un CFL máximo de 0.33 para ese esquema. También advierte dificultades de DG2 con profundidades muy pequeñas/fricción y velocidades espurias que pueden exigir pasos pequeños. Esos datos describen el método publicado en 8.0, no confirman una compilación/binario específico ni opciones de entrada. El artículo trata lluvia como adición al promedio de profundidad; no demuestra soporte de infiltración. No atribuir estos parámetros automáticamente a otra versión.

\prop\ Registrar esquema y versión; comprobar la regla CFL de la implementación concreta, paso adaptativo/fijo, reconstrucción/limitadores, estados seco-húmedo, fronteras y discretización de fricción. Pruebas previas: lago en reposo, conservación de masa, transición seco-húmedo, benchmarks publicados pertinentes, sensibilidad a malla/tiempo y convergencia. El balance debe comparar almacenamiento inicial más aportes menos infiltración y salidas con almacenamiento final, usando volúmenes (m$^3$), intervalos, convenciones y tolerancia aprobada antes de correr. No ocultar inestabilidad recortando resultados sin reportarlo.

\paragraph{Alternativas (no mezclar con la dirección seleccionada).} ATS es aproximación difusiva sin inercia; ACC retiene inercia local pero omite advección; FV1 resuelve SWE completo con esquema de primer orden. Son alternativas identificadas, no partes de DG2. La pertinencia física y capacidades concretas también requieren revisión de versión.

\section{Incertidumbre, validación y límites}
Evaluar incertidumbre de DEM/datum, geometría de cauces naturales, lluvia/caudal, infiltración, rugosidad del terreno, discretización, estado inicial y observación. La exclusión de ciudades no elimina su posible efecto aguas arriba: si altera el flujo, el escenario queda fuera de alcance. Separar calibración de evaluación independiente. Comparar con observación de extensión temporalmente compatible, excluyendo nodata y documentando agua permanente y máscaras. NISAR mide respuesta radar: un cambio no demuestra por sí solo causa, profundidad o velocidad. Si no hay etiquetas adecuadas, el resultado es exploratorio y no se afirma exactitud. Un raster 2D no valida física ni el solver.

\section{Qué se verificó de DG2 y qué no}
La fuente primaria de Shaw et al. (2021), \emph{LISFLOOD-FP 8.0: the new discontinuous Galerkin shallow-water solver for multi-core CPUs and GPUs}, §2.1, ecuaciones (1)--(9), describe ecuaciones y método DG2; §2.1.2 discute el comportamiento en profundidades delgadas/fricción; §2.1.3, ecuación (14), describe tratamiento de lluvia. No se encontró allí soporte de canal subgrid para FV1/DG2. Estas afirmaciones se limitan a lo publicado para 8.0.

Sharifian et al. (2023) estudian LISFLOOD-FP 8.1 y ACC/GPU; las menciones a herramientas/formatos en ese artículo no prueban que una versión concreta de DG2 8.1 admita una determinada opción. Antes de atribuir capacidad a un ejecutable hay que inspeccionar su versión, manual, ayuda, fuente/release y ejecutar una prueba reproducible. No se inspeccionó ni instaló aquí ningún solver. La licencia del artefacto concreto sigue pendiente de verificación; la licencia indicada en publicaciones no resuelve la licencia del binario distribuible.

\part{II. Flujo de aplicación}
\section{Arquitectura y estado: existente frente a propuesto}
\begin{description}
\item[Ramas.] En \path{ETAPA-1} (09cd24d) y \path{simulacion-inundancion} (a5d95fc) existen adquisición e inspección NISAR, además del mismo stack web; la rama remota \path{origin/feat/mapa-editable-sandbox} se observó en 2bc437e. Estas referencias son una instantánea local, no implican que las ramas futuras estén integradas o que ya exista hidráulica.
\item[Frontend.] \path{web/package.json} declara React 19, TypeScript 6, Vite 8, MapLibre GL JS 6 y lucide-react; \path{web/src/components/MapCanvas.tsx} ya tiene mapa, polígonos/ROI, fuentes de imagen y capas raster. \path{web/src/api/client.ts} intenta consultar \path{/api/} pero vuelve a mocks etiquetados si falla. La persistencia local y la simulación solar visual no son resultados hidráulicos. No se requiere otra librería de frontend para una primera vista DG2.
\item[Pipeline/backend.] \path{pipeline/catalog.py}, \path{pipeline/download.py}, \path{pipeline/nisar_h5.py}, \path{pipeline/sources/asf.py}, \path{pipeline/etapa1.py} y \path{pipeline/web_export.py} ya cubren consulta, recorte NISAR e imágenes PNG/COG para exploración; hay un par GCOV T62/F10 recortado, pero no mapa de cambio ni entrada hidráulica derivada. No existen todavía \path{pipeline/hydrology/}, \path{preprocess.py}, \path{change.py} ni \path{validate.py}. \path{backend/config/urls.py} solo registra admin; los modelos de escenas/análisis son esqueletos, no una API de simulaciones.
\item[Propuesta, no instalada.] Ejecutar un solver externo DG2 desde un adaptador de \path{pipeline/hydrology/}, sin dependencia de Django ni proceso hidráulico en el navegador. Django/DRF solo publicaría estado y metadatos verificados; la interfaz existente mostraría resultados diferenciados de observación, mocks y simulación. Los nombres de módulos y contratos hidráulicos de abajo son propuestas, no archivos existentes.
\end{description}

\section{Herramientas para un primer caso no urbano}
\begin{itemize}
\item \textbf{Cálculo hidráulico propuesto:} LISFLOOD-FP 8.0 DG2 en CPU como candidato externo, no paquete Python ni módulo de React. Shaw et al. documentan DG2 para esa versión y la guía de SEAMLESS-WAVE describe compilación y ejemplos; antes de adoptarlo verificar el artefacto, sus instrucciones de ejecución, entradas/salidas, rendimiento y archivo LICENSE. La mención de DG2 en documentación posterior no demuestra soporte concreto de un binario 8.1/8.2/8.3. No instalar ni compilar en esta fase.
\item \textbf{Integración Python:} reutilizar Python del \path{pipeline/}, \texttt{subprocess} y \texttt{json} de la biblioteca estándar para aislamiento, códigos de salida y manifiesto. Reutilizar Rasterio/NumPy ya presentes para validar y alinear DEM/raster numérico, CRS, nodata y COG; \path{web_export.py} es referencia para PNG/COG de observación en dB, no un convertidor hidráulico listo para profundidad en metros. \texttt{h5py} existente lee NISAR como observación, no como DEM ni forzamiento DG2.
\item \textbf{Publicación inicial:} reutilizar React, TypeScript, Vite y MapLibre; una fuente de imagen con esquinas geográficas registradas en el manifiesto y capa raster puede mostrar un PNG para un caso acotado; el COG numérico y el manifiesto, no el PNG, preservan las mediciones. Teselas COG/TiTiler y GPU serían opciones posteriores solo si tamaño, rendimiento y despliegue lo exigen; no agregar dependencias ahora.
\end{itemize}
\textbf{No confundir nombres:} LISFLOOD-FP (hidrodinámica 2D con DG2) no es el proyecto LISFLOOD de hidrología distribuida del JRC. La selección del solver exacto queda pendiente de pruebas y licencia; ninguna librería de visualización sustituye sus ecuaciones.

\section{Recorrido de datos y contratos}
\begin{enumerate}
\item \textbf{Selección de ROI (frontend).} Reutilizar el dibujo de polígonos de MapCanvas/MapLibre; validar polígono, CRS/área y mostrarlo como región de consulta. La ROI no define automáticamente dominio, condición inicial ni frontera hidráulica. Comprobar con datos adecuados si tanto la ROI como el dominio propuesto cumplen el alcance no urbano; si no se puede verificar o hay dependencia hidráulica urbana, informar que este caso simplificado no está disponible. Respuesta: identificador/geom y estado de elegibilidad; errores: geometría inválida, cobertura insuficiente o caso fuera de alcance.
\item \textbf{Disponibilidad y adquisición.} Reutilizar el catálogo y recorte NISAR C1 existentes para observación/contraste, nunca como profundidad simulada. Obtener por separado para la ROI y dominio no urbano propuesto: DEM, cauces naturales, forzamientos, referencia independiente y evidencia de uso del suelo suficiente para evaluar el alcance, con procedencia/licencia. El par T62/F10 disponible es SAR, no DEM ni lluvia. No adquirir capas de edificios, vialidades o drenaje urbano como entradas del modelo. Falta de cobertura o de evidencia para delimitar el dominio bloquea ejecución; no se sustituyen datos faltantes por datos ficticios.
\item \textbf{Preparación.} Con Rasterio/NumPy existentes, determinar dominio computacional, CRS métrico y datum vertical; alinear mallas conservando nodata y método de remuestreo documentado. Separar entradas físicas de escenas radar y de las vistas PNG; no usar un raster de color o dB como elevación. Salida: insumos preparados con metadatos y controles de calidad.
\item \textbf{Configuración estricta.} Declarar alcance no urbano y evidencia de elegibilidad, formulación DG2, solver y versión, periodo/UTC, malla, DEM, fuentes, condiciones iniciales/frontera, infiltración, rugosidad de terreno, umbral de extensión y política de estabilidad. Rechazar áreas urbanas o dependencias hidráulicas urbanas detectadas, campos requeridos ausentes y unidades incoherentes antes de ejecutar; no hay defaults físicos silenciosos.
\item \textbf{Entradas del solver.} Adaptador Python a LISFLOOD-FP DG2 CPU (candidato 8.0) serializa solo formatos y opciones confirmados en el manual y artefacto elegidos, sin implementar un solver SWE propio. No se supone que lluvia, infiltración, ROI, cauces o una opción DG2 se puedan pasar automáticamente.
\item \textbf{Ejecución externa.} Tras verificar licencia y disponibilidad, ejecutar el binario por proceso aislado con \texttt{subprocess} desde \texttt{pipeline/}, sin dependencia de Django, de Vite ni de GPU obligatoria. Registrar versión, argumentos no sensibles, logs, tiempos y código de salida. Timeout, binario/versionado distinto, salida faltante o error numérico invalidan el run; no publicar parciales como resultados válidos.
\item \textbf{Salidas C6 y manifiesto.} Propuesta: \texttt{data/processed/<simulation\_id>/} contiene rasters numéricos de profundidad en metros en tiempos seleccionados, profundidad máxima, extensión con instante/umbral definidos, manifiesto y validación; tiempo de llegada u observación solo si existen y están definidos. Producir PNG de visualización y COG solo tras validar unidades, CRS, tiempos y nodata; no reciclar la escala de dB de NISAR. Manifiesto registra hashes disponibles, fuentes, CRS/transformación/nodata, dominio y ROI, datum, justificación del alcance no urbano, infraestructura excluida, configuración, versiones, log, código, estado y limitaciones; nunca secretos.
\item \textbf{Validación.} Comprobar integridad de rasters, balance, estabilidad, benchmark y referencia independiente. Sin observación adecuada, informar estado exploratorio y límites; no bloquear el registro técnico del run, pero sí cualquier afirmación de validación positiva.
\item \textbf{Backend/API.} Implementación futura sobre Django/DRF ya configurados: persistir estado y metadatos C6/C3 y servir referencias controladas; hoy no hay endpoints hidráulicos. No ejecutar comandos desde petición ni recalcular física. Respuestas distinguen pendiente, fallido, exploratorio y validado; errores de esquema o artefacto corrupto se hacen visibles.
\item \textbf{Visualización y exportación.} Reutilizar MapLibre y las pantallas React existentes para mostrar, tras integrar API y capas verificadas, mapas/series, leyenda, instante, unidades, estado de validación, incertidumbre y procedencia; señalar que solo se modeló terreno no urbano y que no se estiman impactos en ciudades. Mantener separadas simulación y observación. Exportación incluye referencia al manifiesto y al alcance. Sin API/capas válidas, interfaz muestra indisponibilidad, no un resultado simulado.
\end{enumerate}

\section{Criterios de avance y pruebas}
\begin{longtable}{p{.18\linewidth}p{.34\linewidth}p{.36\linewidth}}
\toprule Etapa & Prueba/criterio para avanzar & Fallo observable y tratamiento \\\midrule
ROI y datos & Geometría válida y evidencia del alcance no urbano para ROI y dominio; cobertura, CRS y procedencia inspeccionados. & Sin datos para verificar el alcance, o si hay dependencia urbana: detener el caso e informar qué falta. \\
Dominio/entradas & Alineación, datum, nodata, unidades y configuraciones obligatorias verificados con tests y fixtures. & Huecos, incompatibilidad vertical o campo obligatorio ausente: rechazo previo al solver. \\
Adaptador/solver & Versión exacta documentada; prueba de interfaz; benchmarks DG2, balance, estabilidad y convergencia acordados. & Binario/formato/opción desconocida, timeout o no convergencia: run fallido, logs conservados. \\
C6/manifiesto & Salidas completas y georreferenciadas; manifiesto reproducible, sin secretos; esquema validado. & Archivo faltante, grilla o tiempo incoherente: artefacto no publicable. \\
Validación & Comparación independiente temporal/espacialmente compatible, métricas/máscaras y tolerancias predefinidas. & Sin referencia: etiquetar exploratorio; no inventar exactitud. \\
API, interfaz y exportación & Pruebas de contrato y visualización con C6 de prueba; verificar leyenda, unidades, estado y procedencia. & API o capa ausente: informar indisponibilidad; nunca presentar mocks como simulación real. \\
\bottomrule
\end{longtable}

\section{Decisiones y límites pendientes}
Antes de una implementación física deben confirmarse por responsables: evento/régimen; idoneidad de SWE y DG2 frente a alternativas; elegibilidad no urbana de ROI, dominio y aportes aguas arriba; DEM/datum/cauces; lluvia, caudales y fronteras; infiltración y rugosidad del terreno; versión exacta y formatos de solver; benchmarks, criterios de estabilidad y balance; validación independiente; licencia del artefacto. El primer uso no debe prometer simulación interactiva para cualquier polígono: solo puede ejecutarse donde existan entradas suficientes y configuración revisada. Este documento no implementa simulación, no instala software y no acredita resultados.

\section{Fuentes primarias}
\begin{enumerate}
\item Shaw, J. et al. (2021). ``LISFLOOD-FP 8.0: the new discontinuous Galerkin shallow-water solver for multi-core CPUs and GPUs''. \emph{Geoscientific Model Development}, 14, 3577--3602, §2.1, Eqs. (1)--(9), §2.1.2--2.1.3, Eq. (14). \url{https://gmd.copernicus.org/articles/14/3577/2021/gmd-14-3577-2021.pdf}. DOI: \url{https://doi.org/10.5194/gmd-14-3577-2021}.
\item Sharifian, M. K. et al. (2023). ``LISFLOOD-FP 8.1: new GPU-accelerated solvers for faster fluvial/pluvial flood simulations''. \emph{Geoscientific Model Development}, 16, 2391--2413. \url{https://gmd.copernicus.org/articles/16/2391/2023/}. Referencia de 8.1, no prueba por sí sola capacidades DG2 de un binario.
\item LISFLOOD-FP 8.0, documentación de los desarrolladores SEAMLESS-WAVE (DG2, CPU, compilación y casos): \url{https://www.seamlesswave.com/LISFLOOD8.0}. Fuente y ejemplos del release 8.0: \url{https://zenodo.org/records/4073011}. El archivo LICENSE del artefacto elegido sigue sin verificarse.
\item Rasterio, reproyección y remuestreo de rasters (preparación, no simulación): \url{https://rasterio.readthedocs.io/en/stable/topics/reproject.html}. MapLibre GL JS, fuente de imagen con cuatro esquinas: \url{https://github.com/maplibre/maplibre-gl-js/blob/main/src/source/image_source.ts}.
\item Estado de ramas y contratos: \path{AGENTS.md}, \path{docs/guia-simulacion-inundacion.md}, \path{web/package.json}, \path{web/src/components/MapCanvas.tsx}, \path{web/src/api/client.ts} y \path{pipeline/web_export.py}. Describen código disponible, no resultados de solver.
\end{enumerate}
\end{document}
